\section{Experiments and Results}
\label{sec:experiments_results}

The experimental analysis was designed to follow the central difficulty of this study: multimodal cervical screening data were available at scale, but report-level semantic supervision was sparse and highly uneven across centres. We therefore evaluated LCAD--RASA in a stepwise manner. We first characterised the supervision gap, then tested whether structured pseudo reports could provide modality-grounded semantic anchors rather than simple label-to-text templates. We next examined whether these anchors improved cross-modal alignment and downstream prediction, and finally assessed whether the observed behaviour was stable under ablation, perturbation, and cross-centre evaluation.

\subsection{A five-centre cohort revealed a semantic-supervision bottleneck rather than an imaging-scale bottleneck}
\label{subsec:cohort_results}

The final analytic cohort included 1,897 multimodal cervical screening cases from five centres, with 137,591 evaluable OCT and colposcopy images (Table~\ref{tab:cohort_summary}, ``Cohort scale and case-level report supervision''; Figure~\ref{fig:centre_supervision}, ``Centre-level cohort scale and report-supervision imbalance''). Thus, the primary limitation was not the availability of image-level evidence. Instead, the limiting factor was the availability of structured clinician-authored report supervision. Archived real reports were available for 744 cases, whereas 1,153 cases had complete multimodal evidence but no archived report and were therefore candidates for pseudo-report completion.

This report-supervision imbalance was strongly centre-dependent. Enshi retained real reports for all 406 cases, Jingzhou retained reports for 334 of 406 cases, Xiangyang retained only 4 reports among 500 cases, and Wuhan and Shiyan had no centre-wide real-report archives. This pattern defines the practical learning problem addressed by LCAD--RASA: the model must exploit large-scale multimodal evidence while avoiding the assumption that semantic report supervision is uniformly observed. The subsequent experiments therefore focus on whether missing report supervision can be transformed into auditable semantic alignment signals without treating pseudo reports as equivalent to physician-authored reports.

\subsection{Structured pseudo reports provided modality-grounded semantic anchors rather than label templates}
\label{subsec:pseudo_report_results}

We next examined whether pseudo reports could carry useful modality-specific information, or whether they merely expanded diagnostic labels into repetitive text. The label-template baseline served as a negative control. Although it was schema-complete, it had no modality support, with a mean modality-support rate of 0.000, and showed high textual repetition, with a maximum duplicate fraction of 0.867 (Table~\ref{tab:pseudo_source_comparison}, ``Structured pseudo-report source comparison''; Figure~\ref{fig:pseudo_source_comparison}, ``Structured pseudo-report source comparison across template, rule-based, and local LLM sources''). This result shows that schema completeness alone is insufficient: a generated report can satisfy the output format while still failing to encode OCT, colposcopy, or clinical evidence.

In contrast, the rule-based and local embedding-assisted LLM pseudo reports preserved OCT, colposcopy, and clinical section support, with a mean modality-support rate of 1.000 for both approaches. The rule-based baseline achieved a label-consistency mean of 0.628, and the local LLM baseline achieved the same mean while retaining modality-grounded sections. These findings support the use of structured pseudo reports as semantic anchors for representation learning, rather than as a simple textual restatement of labels.

The external-provider analysis further tested whether pseudo-report generation depended on a single LLM source. Template, rule-based, local embedding LLM, Qwen-Plus, GLM-4.7-Flash, Gemini-3.1-Pro, and GPT-5.5 were evaluated under a common structured-report schema (Table~\ref{tab:llm_provider_comparison}, ``LLM provider comparison for structured pseudo-report generation''; Figure~\ref{fig:api_provider_summary_heatmap}, ``LLM provider comparison for structured pseudo-report generation''; Figure~\ref{fig:api_quality_heatmap}, ``Stage 1 provider quality metrics and risk-oriented quality indicators''; Figure~\ref{fig:api_reliability}, ``Provider latency, modality support, and generation reliability''). All evaluated sources produced schema-valid structured outputs in their available cohorts. GPT-5.5 was evaluated on the extended 100-case cohort and generated 100 valid structured reports, with a schema-valid rate of 1.000, section completeness of 0.810, modality support of 0.810, contradiction rate of 0.050, hallucination rate of 0.000, mean latency of 11.39 seconds per case, and estimated cost of US\$27.73 per 1,000 cases. Qwen-Plus, GLM-4.7-Flash, and Gemini-3.1-Pro were retained as pilot cohorts; their results are therefore interpreted as feasibility evidence rather than definitive head-to-head provider ranking.

Only de-identified structured evidence was submitted to external APIs. Raw images, image paths, patient names, internal patient identifiers, and hospital identifiers were not transmitted. These safeguards preserve the distinction between pseudo-report augmentation as a controllable learning component and physician-authored reporting as a clinical record.

\subsection{Report sections exposed measurable cross-modal semantic alignment}
\label{subsec:alignment_results}

After establishing that pseudo reports could preserve modality-grounded sections, we asked whether these sections were actually aligned with the corresponding multimodal representations. OCT embeddings were matched to \texttt{oct\_findings}, colposcopy embeddings to \texttt{colposcopy\_findings}, clinical embeddings to \texttt{clinical\_context}, and fused representations to \texttt{impression}. Full LCAD--RASA achieved a macro mean reciprocal rank (MRR) of 0.049 across the four sections, whereas removing section alignment reduced macro MRR to 0.022 and simple concatenation fusion remained low at 0.020 (Table~\ref{tab:theme1_alignment_ablation}, ``Direct RASA alignment ablation using modality-section retrieval''; Figure~\ref{fig:alignment_retrieval_mrr}, ``Modality-section retrieval MRR for Theme 1 cross-modal semantic alignment'').

The highest retrieval-only macro MRR was observed for the real-report-only variant, at 0.056. This pattern is informative because it separates semantic retrieval quality from downstream diagnostic utility. Real reports provide strong semantic targets when they exist, but their limited and centre-imbalanced availability restricts their coverage. LCAD--RASA therefore aims not to replace real reports, but to extend report-guided alignment to cases where multimodal evidence is present but report supervision is absent.

The staged API alignment analysis showed additional provider-dependent variation. Gemini-3.1-Pro achieved the highest available pilot macro MRR of 0.366, whereas GPT-5.5 achieved macro MRR of 0.063 on the extended cohort (Figure~\ref{fig:api_alignment_mrr}, ``Stage 2 provider alignment analysis: macro MRR and section-specific MRR''; Figure~\ref{fig:api_section_mrr}, ``Provider-level section MRR across report sections''). Because these evaluations differed in cohort size and stage, the values should not be read as a definitive provider ranking. Instead, they show that structured report sections can be used as explicit alignment targets and that alignment behaviour remains measurable across pseudo-report sources.

\subsection{Pseudo-report augmentation was most beneficial when real reports were scarce}
\label{subsec:scarcity_results}

The next question was whether semantic augmentation remained useful when real-report supervision was deliberately reduced. The report-supervision scarcity curve showed the largest gain in the lowest-supervision regime (Table~\ref{tab:scarcity_curve}, ``Report-supervision scarcity curve using a lightweight downstream surrogate''; Figure~\ref{fig:scarcity_curve}, ``Report-supervision scarcity curve comparing real-report-only and pseudo-report-augmented surrogates''). At 10\% real-report availability, the LCAD-augmented surrogate achieved AUROC 0.818 and F1 0.502, compared with AUROC 0.636 and F1 0.297 for the real-report-only surrogate. At full real-report availability, the augmented surrogate also remained higher, with AUROC 0.841 compared with 0.796 for the real-report-only surrogate.

This trend indicates that pseudo-report augmentation is not merely adding redundant text when physician reports are already abundant. Its main value appears in the clinically realistic regime in which images and structured clinical variables are available, but report-level supervision is incomplete. The staged API downstream surrogate was intentionally conservative: GPT-5.5 achieved AUROC 0.630 and F1 0.229 at 10\% real-report availability on the extended cohort, whereas Gemini-3.1-Pro achieved AUROC 0.725 and F1 0.379 in its pilot cohort (Figure~\ref{fig:api_scarcity_auc}, ``Stage 3 downstream scarcity surrogate for selected API-generated pseudo reports''). These surrogate analyses support the usefulness of structured pseudo-report augmentation under report scarcity, while not replacing full end-to-end retraining evidence.

\subsection{Held-out testing favoured the full LCAD--RASA configuration in point estimates}
\label{subsec:diagnostic_results}

Having verified the semantic-alignment pathway, we evaluated whether the full configuration translated into held-out diagnostic performance. On the locked held-out test set ($n=288$), Full LCAD--RASA achieved AUROC 0.832 (95\% CI 0.757--0.897) and F1 0.611 (95\% CI 0.458--0.737) at the validation-selected threshold of 0.37 (Table~\ref{tab:main_comparison}, ``Held-out test performance at validation-selected thresholds''; Figure~\ref{fig:main_auc}, ``Held-out AUROC with bootstrap confidence intervals''; Figure~\ref{fig:main_metrics}, ``Held-out multi-metric profile across model variants''; Figure~\ref{fig:main_auc_f1}, ``AUROC--F1 trade-off across main model variants''). Pseudo-augmented LCAD training achieved AUROC 0.826 (95\% CI 0.747--0.894) and F1 0.545 (95\% CI 0.395--0.674). Real-report-only training achieved AUROC 0.725, whereas simple concatenation fusion achieved AUROC 0.472.

These results show a consistent ordering of the main configurations: generic concatenation was insufficient, real-report-only training was limited by report coverage, pseudo-augmented training improved over real-report-only learning, and the full LCAD--RASA model achieved the highest point estimates. However, paired bootstrap testing did not support a statistical superiority claim. The valid interpretation is therefore that Full LCAD--RASA produced numerically higher held-out performance under the locked retrospective protocol, rather than definitive evidence of clinical superiority.

\subsection{Ablation results linked the observed gains to section-level alignment and operating-point behaviour}
\label{subsec:ablation_results}

The ablation experiments were then used to determine whether the performance pattern could be attributed to the proposed report-aware alignment design. Removing section alignment substantially reduced direct alignment quality, with macro MRR decreasing from 0.049 for Full LCAD--RASA to 0.022 (Table~\ref{tab:theme1_alignment_ablation}, ``Direct RASA alignment ablation using modality-section retrieval''). In the held-out comparison, the LCAD variant without section alignment retained a similar AUROC of 0.826, but its F1 score decreased to 0.510 compared with 0.611 for Full LCAD--RASA (Table~\ref{tab:main_comparison}, ``Held-out test performance at validation-selected thresholds''). This indicates that section alignment may affect threshold-dependent error balance and decision operating behaviour even when rank-based discrimination remains similar.

Additional modality, quality-control, alignment-weight, and component ablations are reported in Supplementary Tables~S1 and S3--S5 and Figure~\ref{fig:rasa_lambda} (``Alignment-weight sweep for the RASA objective''), Figure~\ref{fig:rasa_component_boxenplot} (``Modality ablation and RASA component ablation''), and Figure~\ref{fig:qc_ablation} (``LCAD quality-control ablation''). Taken together, these ablations support the internal coherence of the RASA design. They should be interpreted as mechanistic diagnostics, not as proof that any single component causally explains the full diagnostic behaviour.

\subsection{Perturbation experiments distinguished modality-grounded decoding from fixed-template generation}
\label{subsec:perturbation_results}

A remaining concern was that generated sections might still follow fixed templates despite satisfying section-level metrics. We therefore perturbed modality evidence and examined whether the corresponding report sections degraded selectively. The upgraded perturbation matrix showed that removal or disruption of a modality produced the largest degradation in the matching section (Table~\ref{tab:perturbation_matrix}, ``Upgraded perturbation sensitivity matrix''; Figure~\ref{fig:perturbation_heatmap}, ``Perturbation condition by report-section similarity matrix''; Figure~\ref{fig:perturbation_lineplot}, ``Section-level degradation profiles under modality perturbations''; Figure~\ref{fig:risk_delta}, ``Risk-score shifts under report-generation perturbations''; Figure~\ref{fig:theme1_perturbation}, ``Theme 1 upgraded perturbation sensitivity matrix''). Masking OCT produced a 0.743 drop in the OCT findings section, masking colposcopy produced a 1.000 drop in the colposcopy findings section, and masking clinical instruction produced a 0.833 drop in the clinical-context section.

Label-only inference produced the largest overall report drop of 0.466 and the largest absolute risk shift of 0.371. This finding is important because it shows that the model's report behaviour was not fully recoverable from class labels alone. The section-specific degradation pattern provides fidelity evidence for modality-grounded decoding and helps distinguish LCAD--RASA from label-copying or schema-filling baselines. The evidence remains perturbational rather than causal; we therefore do not claim that the generated reports constitute causal clinical explanations.

\subsection{Reference-stratified and cross-centre analyses exposed the limits of pooled performance}
\label{subsec:loco_results}

We next asked whether the held-out behaviour depended on whether a case had an archived reference report. Full LCAD--RASA achieved AUROC 0.824 among cases with reference reports ($n=108$) and AUROC 0.850 among cases without reference reports ($n=180$). This suggests that the model's held-out behaviour was not confined to the subset with archived reports, although the analysis remains retrospective and conditioned on the available cohort composition.

Cross-centre evaluation revealed a more restrictive picture. Under strict leave-one-centre-out retraining with the fixed quick-budget protocol, held-out AUROC was 0.702 for Enshi, 0.648 for Jingzhou, and 0.382 for Xiangyang (Supplementary Tables~S2 and S2b; Figure~\ref{fig:loco_heatmap}, ``Leave-one-centre-out evaluation: AUROC heatmap and forest-style centre summary''; Figure~\ref{fig:loco_forest}, ``Strict LOCO centre-wise behaviour under fixed quick-budget retraining''). Eval-only LOCO using a global checkpoint is reported separately and is not pooled with strict LOCO retraining. These results show that centre shift remains a material limitation. Therefore, the evidence supports report-aware fusion under retrospective held-out testing, but not an unrestricted external-validation claim across all centres.

\subsection{Robustness, safety, and scalability analyses supported retrospective feasibility}
\label{subsec:robustness_safety_scalability}

Finally, we examined whether the pipeline was reproducible and computationally feasible at cohort scale. Robustness and safety analyses included masking validation, modality ablations, quality-control ablations, multi-seed stability, report-level safety metrics, and runtime summaries (Supplementary Tables~S7--S11; Supplementary Figure~\ref{fig:supplementary_robustness}, ``Supplementary robustness checks: masking validation and multi-seed stability''). The pipeline processed approximately 137k images, and the publishable model checkpoints were approximately 80 MB, indicating that the workflow is feasible for retrospective large-scale multimodal analytics.

An expert-rating template was prepared, but physician scores were not available for the present manuscript version. Accordingly, we do not make claims of physician validation, expert-validated report quality, clinical deployment readiness, or prospective decision support. The completed experiments support a reproducible retrospective framework for LLM-augmented multimodal fusion under incomplete report supervision. Prospective clinical evaluation and stronger cross-centre validation remain necessary before deployment-oriented claims can be made.
